% ============================================================
% 4. RESULTADOS Y ANÁLISIS
% Organizado por objetivo específico, para que se vea que cada
% uno se cumplió. Cada sección cierra con qué dice el resultado,
% no solo con la figura.
% ============================================================
\chapter{Resultados y análisis}
\label{cap:resultados}

% ------------------------------------------------------------
\section{Verificación de la implementación (objetivo 1)}
\label{sec:res-obj1}

\pendiente{reproducción del código de referencia; error frente a RS y al haz
gaussiano; tiempos en CPU y GPU (resultados/tiempos.png,
resultados/aceleracion.png).}

% ------------------------------------------------------------
\section{Validez bajo el muestreo de DLHM (objetivo 2)}
\label{sec:res-obj2}

\pendiente{mapas de error por método; rangos de distancia, apertura numérica,
$s$ y $K_f$ donde cada uno es confiable; validación de la iluminación
esférica; comparación MPASM--CZT.}

% ------------------------------------------------------------
\section{Reconstrucción de hologramas experimentales (objetivo 3)}
\label{sec:res-obj3}

\pendiente{reconstrucciones con cada propagador frente al ASM con onda
esférica del laboratorio; tabla de métricas; figuras del \emph{target}.}

% ------------------------------------------------------------
\section{Discusión}
\label{sec:discusion}

\pendiente{responder la pregunta de investigación: qué del control de muestreo
mejora la reconstrucción y en qué condiciones de registro conviene cada
método. Incluir los límites de lo que se midió.}
